\section{De los modelos de difusión a los modelos de flujo: una conversión de paradigma}

\subsection{Inversión de la dirección temporal}

Antes de introducir formalmente el \textbf{Flow Matching}, el orador enfatizó un cambio en la convención de notación que es \textbf{extremadamente confuso}:

\begin{warningbox}{Inversión de la dirección temporal — el punto más propenso a confusiones}
La dirección temporal del \textbf{modelo de difusión} (desde la perspectiva del proceso hacia adelante):
\begin{itemize}
    \item $t = 0$: distribución de datos $p_{\text{data}}$
    \item $t = T$: distribución de ruido (gaussiana estándar) $\mathcal{N}(0, I)$
\end{itemize}

La dirección temporal del \textbf{Flow Matching} (desde la perspectiva del proceso generativo):
\begin{itemize}
    \item $t = 0$: distribución de ruido (gaussiana estándar) $p_0 = \mathcal{N}(0, I)$
    \item $t = 1$: distribución de datos $p_1 = p_{\text{data}}$
\end{itemize}

Es decir, en el Flow Matching, $x_0$ representa una muestra de ruido gaussiano y $x_1$ representa una muestra de datos reales. ¡Esto es completamente opuesto al hábito establecido en los modelos de difusión!
\end{warningbox}

El orador reconoció que incluso él mismo se confunde frecuentemente entre estas dos convenciones de notación. La razón por la que se mantiene esta inconsistencia es que los modelos de difusión y el flujo matching provienen de tradiciones literarias distintas; sus respectivos artículos e libros de texto siguen sus propias convenciones.

\subsection{Diferencias clave entre modelos de difusión y modelos de flujo}

\begin{importantbox}{Diferencia central: flexibilidad de la distribución de referencia}
En los \textbf{modelos de difusión}, todas las deducciones asumen que la distribución de referencia (distribución de ruido) es una gaussiana estándar $\mathcal{N}(0, I)$. Si se cambiara a otra distribución (por ejemplo, una gaussiana anisotrópica), todas las deducciones matemáticas tendrían que realizarse de nuevo.

Una gran ventaja de los \textbf{modelos de flujo} es que la distribución de referencia $p_0$ puede ser \textbf{cualquier distribución}. Sin importar cómo cambie $p_0$, las fórmulas matemáticas del flujo matching permanecen inalteradas. Esto permite considerar a los modelos de flujo como una \textbf{generalización} de los modelos de difusión.
\end{importantbox}

Aunque ambos provienen de tradiciones literarias distintas, finalmente son matemáticamente equivalentes (en el caso especial de una distribución de referencia gaussiana estándar). Este es un hallazgo importante que esta clase busca revelar.

\subsection{Resumen de la sección}

El Flow Matching adopta la perspectiva del proceso generativo, invierte la dirección temporal y permite cualquier distribución de referencia. Aunque finalmente se puede establecer una relación de equivalencia con los modelos de difusión, esta nueva perspectiva aporta objetivos de entrenamiento más concisos y un espacio de diseño de modelos más flexible.
